\section{Methodology}

This section describes the design principles, research methodology, and evaluation framework guiding Malak Platform's development.

\subsection{Design Principles}

\subsubsection{End-to-End Optimization}
Rather than treating training, compression, compilation, and deployment as independent stages, we co-design these phases to maximize overall system efficiency. For instance, quantization-aware training anticipates deployment constraints, while EdgeIR optimization passes exploit knowledge of target hardware characteristics.

\subsubsection{Hardware-Software Co-Design}
The platform abstracts hardware diversity through a layered architecture:
\begin{enumerate}
\item \textbf{Operator interface}: High-level tensor operations independent of backend
\item \textbf{Backend plugins}: Hardware-specific implementations (ARM NEON, CUDA, NPU)
\item \textbf{Autotuning}: Per-device optimization discovering optimal kernel configurations
\end{enumerate}

This approach balances portability with performance, enabling drop-in hardware upgrades without application rewrites.

\subsubsection{Production-Aware Development}
From inception, we prioritize features critical for real-world deployment:
\begin{itemize}
\item \textbf{Deterministic behavior}: Fixed-point arithmetic eliminates floating-point nondeterminism
\item \textbf{Graceful degradation}: Fallback paths handle out-of-distribution inputs
\item \textbf{Auditability}: Telemetry logs enable post-deployment debugging
\item \textbf{Privacy by design}: On-device processing minimizes data exposure
\end{itemize}

\subsubsection{Developer Ergonomics}
Unified CLI (\texttt{edgeai}) and declarative model cards reduce cognitive load, allowing developers to focus on application logic rather than toolchain intricacies. YAML configurations specify high-level intent, with the framework handling low-level details.

\subsection{Research Methodology}

\subsubsection{Application-Driven Validation}
We validate Malak Platform through three diverse applications representing distinct challenge domains:

\begin{enumerate}
\item \textbf{Vision Jellyfish}: Real-time marine life detection testing vision pipeline optimization, data augmentation strategies, and outdoor deployment robustness
\item \textbf{Energy Optimization}: Time-series forecasting for smart grids evaluating temporal model compression and sequential data handling
\item \textbf{Neuro MS}: Medical image segmentation assessing privacy-preserving inference, high-resolution processing, and clinical accuracy requirements
\end{enumerate}

These applications stress different platform components: convolutional networks (vision), recurrent architectures (energy), attention mechanisms (medical), and mixed data modalities.

\subsubsection{Benchmarking Methodology}
We adopt a standardized evaluation protocol:

\textbf{Baseline Models}: Train full-precision (FP32) models using standard PyTorch training procedures until convergence on validation sets.

\textbf{Compression Pipeline}: Apply quantization (INT8, INT4), pruning (30\%, 50\%, 70\% sparsity), and knowledge distillation, measuring accuracy impact at each stage.

\textbf{Compilation}: Generate optimized binaries for target devices (Cortex-M7 @ 480MHz, Raspberry Pi 4, NVIDIA Jetson Nano) using TVM, MLIR, or XLA backends.

\textbf{Profiling}: Measure per-layer latency, end-to-end throughput, memory consumption (peak RAM, flash size), and energy per inference using on-chip performance counters and external power monitors.

\textbf{Accuracy Validation}: Evaluate compressed models on held-out test sets, computing task-specific metrics (mAP for detection, MAE for regression, Dice coefficient for segmentation).

\subsection{Evaluation Metrics}

\subsubsection{Performance Metrics}
\begin{itemize}
\item \textbf{Latency}: Per-inference time (milliseconds) from input to output
\item \textbf{Throughput}: Inferences per second under sustained load
\item \textbf{Energy}: Joules per inference measured at device power rails
\item \textbf{Memory}: Peak RAM usage and flash storage requirements
\end{itemize}

\subsubsection{Accuracy Metrics}
Task-dependent metrics:
\begin{itemize}
\item \textbf{Vision}: Mean Average Precision (mAP), classification accuracy
\item \textbf{Energy}: Mean Absolute Error (MAE), Root Mean Squared Error (RMSE)
\item \textbf{Medical}: Dice coefficient, Intersection over Union (IoU)
\end{itemize}

\subsubsection{Compression Efficiency}
\begin{itemize}
\item \textbf{Compression ratio}: Model size reduction (compressed / baseline)
\item \textbf{Accuracy degradation}: Percentage point loss vs. FP32 baseline
\item \textbf{Pareto efficiency}: Latency-accuracy trade-off frontier
\end{itemize}

\subsubsection{Production Readiness}
\begin{itemize}
\item \textbf{Telemetry overhead}: Runtime cost of monitoring instrumentation
\item \textbf{Drift detection sensitivity}: False positive/negative rates for distribution shifts
\item \textbf{Privacy compliance}: Formal verification of on-device data policies
\end{itemize}

\subsection{Hardware Testbeds}

Table \ref{tab:testbeds} lists target devices spanning microcontrollers to edge servers:

\begin{table}[h]
\centering
\caption{Hardware Testbeds for Evaluation}
\label{tab:testbeds}
\resizebox{\columnwidth}{!}{%
\begin{tabular}{lcccc}
\hline
\textbf{Device} & \textbf{CPU} & \textbf{Clock} & \textbf{RAM} & \textbf{Accelerator} \\
\hline
STM32H7 & Cortex-M7 & 480 MHz & 1 MB & - \\
Raspberry Pi 4 & Cortex-A72 & 1.5 GHz & 4 GB & - \\
NVIDIA Jetson Nano & Cortex-A57 & 1.43 GHz & 4 GB & 128-core GPU \\
ESP32-S3 & Xtensa LX7 & 240 MHz & 512 KB & - \\
Coral Dev Board & Cortex-A53 & 1.5 GHz & 1 GB & Edge TPU \\
\hline
\end{tabular}%
}
\end{table}

This diversity validates portability across ARM Cortex-M (ultra-low-power), Cortex-A (mid-range), and accelerated platforms (GPU, TPU).

\subsection{Open Science Practices}

To ensure reproducibility:
\begin{itemize}
\item \textbf{Public repository}: All code, model cards, and benchmarks available at \url{https://github.com/alnemari-m/malak_platform}
\item \textbf{Docker images}: Pre-configured environments with pinned dependencies
\item \textbf{Benchmark datasets}: Training/validation splits with preprocessing scripts
\item \textbf{Automated CI}: Continuous integration tests validate cross-platform builds
\end{itemize}

Researchers can replicate experiments via:
\begin{verbatim}
git clone https://github.com/alnemari-m/malak_platform
cd malak_platform/malak_edge_ai
./reproduce_experiments.sh
\end{verbatim}
